\chapter{地址映射与指令编码}
\renewcommand{\thesubsection}{\thechapter.\arabic{subsection}}

\Needspace{150pt}

\section{SoC 地址映射}\label{soc-ux5730ux5740ux6620ux5c04}

{\small
\setlength{\tabcolsep}{4pt}
\renewcommand{\thetable}{B-1}
\begin{longtable}{@{}>{\raggedright\arraybackslash}p{\dimexpr 0.1496\linewidth-2\tabcolsep\relax}>{\raggedright\arraybackslash}p{\dimexpr 0.2430\linewidth-2\tabcolsep\relax}>{\raggedright\arraybackslash}p{\dimexpr 0.6073\linewidth-2\tabcolsep\relax}@{}}
\caption{SoC 地址映射}\label{tab:B-1}\\
\toprule
\textbf{窗口} & \textbf{基址} & \textbf{从设备} \\
\midrule
\endfirsthead
\multicolumn{3}{@{}l}{\footnotesize 续表 B-1}\\
\toprule
\textbf{窗口} & \textbf{基址} & \textbf{从设备} \\
\midrule
\endhead
MEM & 0x10000000 & 128 KiB AXI BRAM，统一代码数据空间 \\
CLINT & 0x12000000 & 定时器与软件中断 \\
UART & 0x20000000 & 当前 mailbox 的 RV32 uncached MMIO 端 \\
GPIO & 0x20010000 & LED 与按键事件 \\
ASCON & 0x30000000 & src/rtl/ascon/ascon\_mmio 的 AEAD 寄存器 \\
PUF KEY & 0x30000080 & PUF3 变体 key\_mmio 的状态、控制与诊断回读 \\
CPI & 0x60000000 & 指令与停顿等事件，不能相加替代总周期 \\
\bottomrule
\end{longtable}
}

\Needspace{85pt}

\section{定制指令编码}

主操作码 custom-0=0x0B、funct7=0；funct3=0/1/2/3 分别为 ascon.w/ascon.rd/ascon.p12/ascon.p8。状态共 10 个 32-bit 字，按 s0 低/高至 s4 低/高排列；非法索引读零、写忽略。索引取自寄存器值而非寄存器编号；具体接口与取消边界见 §2.3.3.4，不声称完整非法编码空间已验收。

\Needspace{85pt}

\section{MMIO 加速器寄存器}\label{mmio-ux52a0ux901fux5668ux5bc4ux5b58ux5668}

AEAD 寄存器为 CTRL(0x00)、STATUS(0x04)、ADLEN(0x08)、MLEN(0x0C)、KEY(0x10--1C)、NONCE(0x20--2C)、TAGIN(0x30--3C)、AD\_WDATA(0x40)、MSG\_WDATA(0x44)、OUT\_RDATA(0x48)、TAGOUT(0x4C--58)。标准 KAT/benchmark 使用调用者 key；HWKEY 开启时忽略 KEY 写入，START 锁存内部派生 key。

PUF3 在同一基址的 +0x80 为 STATUS，+0x84 为 CTRL（[0] START、[1] ENROLL、[2] CLEAR、[3] HWKEY\_ON、[4] HWKEY\_OFF，OFF 优先）；STATUS[6]/[7] 为实际选择/开关请求。+0x88/8C 为 DNA、+0x90 为 QUALITY、+0x94--A0 为 REF、+0xA4--B0 为 MASK、+0xB4--C0 为 DITHER、+0xC4--D0 为 RAW。这些诊断量可读，key 无直接读地址不等于保密隔离。

\Needspace{150pt}

\section{PYNQ AXI-Lite 窗口}\label{pynq-axi-lite-ux7a97ux53e3}

{\small
\setlength{\tabcolsep}{4pt}
\renewcommand{\thetable}{B-2}
\begin{longtable}{@{}>{\raggedright\arraybackslash}p{\dimexpr 0.1548\linewidth-2\tabcolsep\relax}>{\raggedright\arraybackslash}p{\dimexpr 0.3115\linewidth-2\tabcolsep\relax}>{\raggedright\arraybackslash}p{\dimexpr 0.5337\linewidth-2\tabcolsep\relax}@{}}
\caption{PYNQ AXI-Lite 窗口}\label{tab:B-2}\\
\toprule
\textbf{窗口} & \textbf{PS 基址} & \textbf{功能} \\
\midrule
\endfirsthead
\multicolumn{3}{@{}l}{\footnotesize 续表 B-2}\\
\toprule
\textbf{窗口} & \textbf{PS 基址} & \textbf{功能} \\
\midrule
\endhead
CONTROL & 0x40000000（4 KiB） & 复位、按键控制、LED/复位状态；ID=0x33465150（PQF3）、VERSION=1 \\
MAILBOX & 0x40001000（4 KiB） & 字节 TX/RX FIFO、状态、overflow、计数、ID \\
\bottomrule
\end{longtable}
}

地址由当前 HWH 校验；不是 RV32 本地 MMIO 地址。详细定义见 §2.3.3.7。
